\documentclass[12pt,a4paper]{article}

\usepackage[T2A]{fontenc}
\usepackage[utf8]{inputenc}
\usepackage[russian]{babel}
\usepackage{amsmath,amssymb,amsthm}
\usepackage{enumitem}
\usepackage{array}
\usepackage{longtable}
\usepackage{geometry}
\usepackage{tikz}
\usepackage{booktabs}
\usepackage{hyperref}

\geometry{margin=2.2cm}

\newtheorem{definition}{Определение}
\newtheorem{remark}{Замечание}
\newtheorem{example}{Пример}
\newtheorem{principle}{Принцип}

\title{Билет №60 \\ \small Управление проектом, понятие проекта. Творческие и индустриальные проекты. Типичные фазы индустриального проекта. Сбор требований, планирование. План. (СпБГУ)}
\author{Сериков Владислав Александрович}
\date{16 мая 2026}

\begin{document}

\maketitle
\tableofcontents
\newpage

\section{Понятие проекта}

\begin{definition}
\textbf{Проект} --- это временная целенаправленная деятельность, выполняемая в условиях ограничений и неопределённости для создания уникального результата: продукта, услуги, процесса, организационного изменения или иного ценного результата.
\end{definition}

Ключевые признаки проекта:
\begin{enumerate}
    \item \textbf{Временность}: проект имеет начало и завершение.
    \item \textbf{Уникальность результата}: результат проекта отличается от полностью повторяемой операционной деятельности.
    \item \textbf{Цель}: проект существует ради достижения заранее сформулированной цели или набора целей.
    \item \textbf{Ограничения}: проект ограничен сроками, бюджетом, ресурсами, качеством, технологиями, нормативными требованиями.
    \item \textbf{Неопределённость}: в проекте всегда присутствуют риски, изменения требований, технические и организационные неопределённости.
    \item \textbf{Заинтересованные стороны}: на проект влияют заказчик, пользователи, команда, руководство, поставщики, регулирующие органы и другие участники.
\end{enumerate}

\begin{definition}
\textbf{Управление проектом} --- это применение знаний, методов, инструментов и организационных практик к проектной деятельности для достижения целей проекта с учётом ограничений по содержанию, срокам, стоимости, качеству, ресурсам и рискам.
\end{definition}

\section{Проект и операционная деятельность}

Проектную деятельность важно отличать от операционной.

\begin{longtable}{|p{0.28\textwidth}|p{0.33\textwidth}|p{0.33\textwidth}|}
\hline
\textbf{Критерий} & \textbf{Проект} & \textbf{Операционная деятельность} \\
\hline
Цель & Создать уникальный результат & Поддерживать регулярное функционирование организации \\
\hline
Временные рамки & Есть начало и конец & Непрерывна или циклична \\
\hline
Повторяемость & Низкая или средняя & Высокая \\
\hline
Неопределённость & Обычно высокая & Обычно ниже, процессы стандартизованы \\
\hline
Пример & Разработка новой информационной системы & Ежедневная техническая поддержка пользователей \\
\hline
\end{longtable}

\begin{example}
Разработка нового мобильного приложения --- проект, потому что имеет уникальный результат, ограничения по срокам и бюджету. Поддержка уже выпущенного приложения --- операционная деятельность, хотя отдельный крупный релиз может оформляться как проект.
\end{example}

\section{Основные ограничения проекта}

Классически выделяют так называемый \textbf{треугольник проекта}:

\[
\text{Содержание} \quad \leftrightarrow \quad \text{Сроки} \quad \leftrightarrow \quad \text{Стоимость}.
\]

Позднее модель расширяют, добавляя:
\begin{itemize}
    \item качество;
    \item риски;
    \item ресурсы;
    \item удовлетворённость заинтересованных сторон;
    \item ценность для бизнеса или общества.
\end{itemize}

\begin{center}
\begin{tikzpicture}[scale=1.1]
    \draw[thick] (0,0) -- (4,0) -- (2,3.2) -- cycle;
    \node at (2,-0.35) {Стоимость};
    \node at (-0.35,1.5) {Сроки};
    \node at (4.35,1.5) {Содержание};
    \node at (2,1.3) {Качество};
\end{tikzpicture}
\end{center}

Если увеличить содержание проекта без изменения сроков и бюджета, то обычно страдает качество или возрастает риск. Если сократить сроки, может потребоваться больше ресурсов или уменьшение объёма работ.

\section{Творческие и индустриальные проекты}

\subsection{Творческие проекты}

\begin{definition}
\textbf{Творческий проект} --- это проект, в котором существенная часть результата заранее не может быть полностью формализована, а успех зависит от поиска новой идеи, художественного, научного, исследовательского или инженерного решения.
\end{definition}

Примеры:
\begin{itemize}
    \item создание фильма, спектакля, игры, музыкального альбома;
    \item разработка концепции нового продукта;
    \item научно-исследовательская работа;
    \item экспериментальный стартап;
    \item проектирование принципиально нового пользовательского опыта.
\end{itemize}

Характерные признаки творческих проектов:
\begin{enumerate}
    \item высокая неопределённость результата;
    \item сложность точной оценки сроков и стоимости на ранней стадии;
    \item необходимость итераций, проб и ошибок;
    \item высокая роль экспертизы и креативности команды;
    \item возможность изменения цели по мере получения нового знания;
    \item важность прототипов, демонстраций, пользовательской обратной связи.
\end{enumerate}

Для творческих проектов часто применяют гибкие и итерационные подходы:
\begin{itemize}
    \item Agile;
    \item Scrum;
    \item Kanban;
    \item Lean Startup;
    \item дизайн-мышление;
    \item прототипирование.
\end{itemize}

\subsection{Индустриальные проекты}

\begin{definition}
\textbf{Индустриальный проект} --- это проект, направленный на создание материального или технологически сложного результата по заранее определённым требованиям, стандартам, регламентам и производственным ограничениям.
\end{definition}

Примеры:
\begin{itemize}
    \item строительство завода, моста, линии метро;
    \item внедрение производственной линии;
    \item разработка серийного промышленного изделия;
    \item внедрение корпоративной информационной системы;
    \item создание инфраструктурного объекта;
    \item сертификация и запуск медицинского или авиационного оборудования.
\end{itemize}

Характерные признаки индустриальных проектов:
\begin{enumerate}
    \item высокая стоимость ошибки;
    \item значительная роль нормативов, стандартов и документации;
    \item необходимость детального планирования;
    \item формальные процедуры контроля качества;
    \item большое число зависимостей между работами;
    \item важность управления поставками, контрактами и рисками;
    \item наличие стадийных проверок и утверждений.
\end{enumerate}

\subsection{Сравнение творческих и индустриальных проектов}

\begin{longtable}{|p{0.25\textwidth}|p{0.35\textwidth}|p{0.35\textwidth}|}
\hline
\textbf{Критерий} & \textbf{Творческий проект} & \textbf{Индустриальный проект} \\
\hline
Результат & Часто уточняется в ходе работы & Обычно задаётся заранее через требования и спецификации \\
\hline
Неопределённость & Высокая & Средняя или высокая, но формально управляемая \\
\hline
Метод управления & Итерационный, гибкий, исследовательский & Плановый, стадийный, регламентированный \\
\hline
Документация & Может быть облегчённой & Обычно обширная и формальная \\
\hline
Изменения & Часто естественны и ожидаемы & Управляются через процедуры контроля изменений \\
\hline
Критерии успеха & Новизна, ценность, принятие аудиторией или пользователями & Соответствие требованиям, срокам, бюджету, качеству, нормативам \\
\hline
\end{longtable}

\section{Жизненный цикл проекта}

\begin{definition}
\textbf{Жизненный цикл проекта} --- это совокупность фаз, через которые проходит проект от инициации до завершения.
\end{definition}

Фаза проекта --- логически завершённый этап, после которого обычно принимается управленческое решение:
\begin{itemize}
    \item продолжать проект;
    \item изменить план;
    \item вернуться к предыдущему этапу;
    \item приостановить проект;
    \item закрыть проект.
\end{itemize}

Такие контрольные точки называют \textbf{фазовыми воротами} или \textbf{stage-gate}.

\section{Типичные фазы индустриального проекта}

Для индустриального проекта характерна стадийная модель. В общем виде можно выделить следующие фазы:

\[
\text{Инициация}
\rightarrow
\text{Сбор требований}
\rightarrow
\text{Планирование}
\rightarrow
\text{Проектирование}
\rightarrow
\text{Реализация}
\rightarrow
\text{Тестирование и приёмка}
\rightarrow
\text{Ввод в эксплуатацию}
\rightarrow
\text{Закрытие}.
\]

\subsection{Фаза 1. Инициация}

Цель фазы --- обосновать необходимость проекта и принять решение о его запуске.

Основные действия:
\begin{enumerate}
    \item выявление проблемы или возможности;
    \item определение предварительной цели проекта;
    \item анализ заинтересованных сторон;
    \item предварительная оценка выгод, затрат и рисков;
    \item назначение руководителя проекта;
    \item подготовка устава проекта.
\end{enumerate}

Типовые результаты:
\begin{itemize}
    \item бизнес-обоснование;
    \item устав проекта;
    \item предварительное описание содержания;
    \item список ключевых заинтересованных сторон;
    \item решение о запуске или отказе от проекта.
\end{itemize}

\begin{definition}
\textbf{Устав проекта} --- это документ, формально санкционирующий существование проекта и предоставляющий руководителю проекта полномочия использовать ресурсы организации.
\end{definition}

\subsection{Фаза 2. Сбор и анализ требований}

Цель фазы --- понять, какой результат должен быть создан и каким ограничениям он должен соответствовать.

Основные действия:
\begin{enumerate}
    \item идентификация заинтересованных сторон;
    \item выявление потребностей заказчика и пользователей;
    \item сбор функциональных, нефункциональных, технических, нормативных и бизнес-требований;
    \item анализ противоречий;
    \item приоритизация требований;
    \item формализация требований;
    \item согласование требований.
\end{enumerate}

Типовые результаты:
\begin{itemize}
    \item спецификация требований;
    \item матрица трассируемости требований;
    \item список допущений и ограничений;
    \item критерии приёмки результата.
\end{itemize}

\subsection{Фаза 3. Планирование}

Цель фазы --- определить, как именно проект будет выполнен.

Планирование включает:
\begin{itemize}
    \item определение содержания проекта;
    \item декомпозицию работ;
    \item оценку длительности и стоимости;
    \item построение календарного графика;
    \item планирование ресурсов;
    \item планирование качества;
    \item планирование рисков;
    \item планирование коммуникаций;
    \item планирование закупок;
    \item планирование управления изменениями.
\end{itemize}

Типовой результат --- \textbf{план управления проектом}.

\subsection{Фаза 4. Проектирование}

Цель фазы --- разработать техническое решение, которое удовлетворяет требованиям.

Основные действия:
\begin{enumerate}
    \item выбор архитектуры решения;
    \item разработка проектной документации;
    \item моделирование процессов и объектов;
    \item инженерные расчёты;
    \item выбор технологий, материалов, оборудования;
    \item экспертиза проектных решений;
    \item согласование технической документации.
\end{enumerate}

Типовые результаты:
\begin{itemize}
    \item технический проект;
    \item рабочая документация;
    \item архитектурные схемы;
    \item спецификации оборудования;
    \item проектные расчёты.
\end{itemize}

\subsection{Фаза 5. Реализация}

Цель фазы --- создать результат проекта.

Основные действия:
\begin{itemize}
    \item производство работ;
    \item разработка программного или аппаратного компонента;
    \item закупка и поставка материалов;
    \item монтаж оборудования;
    \item координация подрядчиков;
    \item контроль выполнения графика;
    \item управление качеством;
    \item управление изменениями.
\end{itemize}

\subsection{Фаза 6. Тестирование, контроль качества и приёмка}

Цель фазы --- убедиться, что результат соответствует требованиям.

Основные действия:
\begin{enumerate}
    \item разработка программы испытаний;
    \item проведение тестов и проверок;
    \item выявление дефектов;
    \item устранение несоответствий;
    \item повторные проверки;
    \item оформление актов приёмки.
\end{enumerate}

Типовые результаты:
\begin{itemize}
    \item протоколы испытаний;
    \item отчёты о дефектах;
    \item акты устранения замечаний;
    \item акт приёмки результата.
\end{itemize}

\subsection{Фаза 7. Ввод в эксплуатацию}

Цель фазы --- передать результат в реальное использование.

Основные действия:
\begin{itemize}
    \item обучение пользователей и персонала;
    \item миграция данных, если проект информационный;
    \item пусконаладочные работы;
    \item опытная эксплуатация;
    \item подготовка эксплуатационной документации;
    \item передача результата заказчику.
\end{itemize}

\subsection{Фаза 8. Закрытие проекта}

Цель фазы --- формально завершить проект.

Основные действия:
\begin{enumerate}
    \item подтверждение выполнения обязательств;
    \item закрытие контрактов;
    \item финансовое закрытие;
    \item архивирование документации;
    \item анализ достигнутых результатов;
    \item фиксация извлечённых уроков;
    \item освобождение ресурсов команды.
\end{enumerate}

\section{Сбор требований}

\subsection{Понятие требования}

\begin{definition}
\textbf{Требование} --- это документированное условие, возможность, характеристика или ограничение, которому должен соответствовать результат проекта или процесс его создания.
\end{definition}

Требование отвечает на вопрос: \textit{что должно быть достигнуто или обеспечено?}

\subsection{Классификация требований}

\subsubsection{Бизнес-требования}

Бизнес-требования описывают, какую проблему решает проект и какую ценность он создаёт.

\begin{example}
Сократить время обработки заявки клиента с 5 дней до 1 дня.
\end{example}

\subsubsection{Пользовательские требования}

Пользовательские требования описывают потребности конечных пользователей.

\begin{example}
Оператор должен иметь возможность видеть список всех заявок, ожидающих обработки.
\end{example}

\subsubsection{Функциональные требования}

Функциональные требования описывают функции системы или продукта.

\begin{example}
Система должна автоматически отправлять уведомление клиенту после изменения статуса заявки.
\end{example}

\subsubsection{Нефункциональные требования}

Нефункциональные требования описывают свойства качества:
\begin{itemize}
    \item производительность;
    \item надёжность;
    \item безопасность;
    \item удобство использования;
    \item масштабируемость;
    \item сопровождаемость;
    \item совместимость.
\end{itemize}

\begin{example}
Система должна обрабатывать не менее 1000 запросов в минуту при среднем времени ответа не более 2 секунд.
\end{example}

\subsubsection{Ограничения}

Ограничения задают рамки проекта.

\begin{example}
Решение должно быть реализовано на существующей инфраструктуре заказчика.
\end{example}

\subsubsection{Критерии приёмки}

Критерии приёмки определяют, при каких условиях результат считается принятым.

\begin{example}
Функция считается реализованной, если пользователь может создать заявку, сохранить её, изменить статус и найти её через поиск.
\end{example}

\section{Методы сбора требований}

Основные методы:
\begin{enumerate}
    \item интервью;
    \item анкетирование;
    \item наблюдение за работой пользователей;
    \item анализ документов;
    \item семинары и рабочие сессии;
    \item мозговой штурм;
    \item прототипирование;
    \item анализ аналогичных систем;
    \item моделирование бизнес-процессов;
    \item изучение нормативных требований.
\end{enumerate}

\subsection{Интервью}

Интервью применяется, когда нужно получить детальное понимание потребностей конкретных заинтересованных сторон.

Преимущества:
\begin{itemize}
    \item глубина информации;
    \item возможность задавать уточняющие вопросы;
    \item выявление скрытых проблем.
\end{itemize}

Недостатки:
\begin{itemize}
    \item субъективность;
    \item большие затраты времени;
    \item риск неполноты информации.
\end{itemize}

\subsection{Наблюдение}

Наблюдение полезно, когда пользователи не могут точно сформулировать свои действия или реальные процессы отличаются от регламентов.

\begin{example}
При автоматизации склада аналитик наблюдает, как кладовщик реально принимает товар, где возникают задержки и какие операции выполняются неформально.
\end{example}

\subsection{Прототипирование}

Прототипирование используется, когда требования трудно сформулировать абстрактно.

\begin{example}
Для будущей информационной системы создаётся макет интерфейса. Пользователь смотрит на экран и уточняет: какие поля нужны, какие кнопки лишние, какие действия должны выполняться быстрее.
\end{example}

\section{Алгоритм сбора требований}

\subsection{Общая схема}

\begin{enumerate}
    \item Определить цели проекта.
    \item Выявить заинтересованные стороны.
    \item Выбрать методы сбора требований.
    \item Провести сбор информации.
    \item Зафиксировать требования.
    \item Проверить требования на качество.
    \item Разрешить противоречия.
    \item Приоритизировать требования.
    \item Согласовать требования.
    \item Организовать управление изменениями требований.
\end{enumerate}

\subsection{Иллюстрация алгоритма}

\begin{center}
\begin{tikzpicture}[
    node distance=1.1cm,
    every node/.style={draw, rounded corners, align=center, minimum width=6cm, minimum height=0.75cm},
    arrow/.style={->, thick}
]
\node (a) {Цели проекта};
\node (b) [below of=a] {Заинтересованные стороны};
\node (c) [below of=b] {Сбор информации};
\node (d) [below of=c] {Документирование требований};
\node (e) [below of=d] {Проверка и согласование};
\node (f) [below of=e] {Базовая версия требований};

\draw[arrow] (a) -- (b);
\draw[arrow] (b) -- (c);
\draw[arrow] (c) -- (d);
\draw[arrow] (d) -- (e);
\draw[arrow] (e) -- (f);
\end{tikzpicture}
\end{center}

\subsection{Минимальный пример}

Пусть проект --- разработка системы записи студентов на консультации.

\begin{enumerate}
    \item \textbf{Цель}: уменьшить число конфликтов расписания и ручных согласований.
    \item \textbf{Заинтересованные стороны}: студенты, преподаватели, учебный офис.
    \item \textbf{Сбор}: интервью со студентами и преподавателями.
    \item \textbf{Требования}:
    \begin{itemize}
        \item студент должен видеть свободные слоты преподавателя;
        \item студент должен записаться на свободный слот;
        \item преподаватель должен отменить слот;
        \item система должна отправлять уведомление об изменении записи;
        \item время ответа интерфейса не должно превышать 2 секунд при обычной нагрузке.
    \end{itemize}
    \item \textbf{Приоритеты}:
    \begin{itemize}
        \item запись на слот --- высокий;
        \item уведомления --- высокий;
        \item экспорт расписания в календарь --- средний;
        \item выбор темы оформления интерфейса --- низкий.
    \end{itemize}
    \item \textbf{Согласование}: заказчик утверждает перечень требований как базовую версию.
\end{enumerate}

\section{Качество требований}

Хорошее требование должно быть:
\begin{enumerate}
    \item \textbf{необходимым}: без него цель проекта не достигается;
    \item \textbf{однозначным}: допускает только одну интерпретацию;
    \item \textbf{проверяемым}: можно установить, выполнено оно или нет;
    \item \textbf{реализуемым}: может быть выполнено в рамках ограничений;
    \item \textbf{согласованным}: не противоречит другим требованиям;
    \item \textbf{полным}: содержит всю существенную информацию;
    \item \textbf{трассируемым}: можно проследить связь с целью, задачей, тестом или элементом продукта.
\end{enumerate}

\begin{example}
Плохое требование: ``Система должна быть быстрой''.

Хорошее требование: ``При одновременной работе 500 пользователей 95\% запросов поиска должны обрабатываться не более чем за 2 секунды''.
\end{example}

\section{Матрица трассируемости требований}

\begin{definition}
\textbf{Матрица трассируемости требований} --- это таблица, связывающая требования с бизнес-целями, элементами проекта, задачами реализации, тестами и критериями приёмки.
\end{definition}

Пример:

\begin{center}
\begin{tabular}{|c|p{4cm}|p{3cm}|p{3cm}|p{3cm}|}
\hline
ID & Требование & Бизнес-цель & Компонент & Тест \\
\hline
R1 & Студент записывается на свободный слот & Снизить ручные согласования & Модуль записи & T1 \\
\hline
R2 & Преподаватель отменяет слот & Управлять расписанием & Модуль расписания & T2 \\
\hline
R3 & Система отправляет уведомление & Предотвратить пропуски & Модуль уведомлений & T3 \\
\hline
\end{tabular}
\end{center}

Польза матрицы:
\begin{itemize}
    \item помогает контролировать полноту реализации;
    \item облегчает анализ влияния изменений;
    \item связывает требования с тестированием;
    \item снижает риск потери требований.
\end{itemize}

\section{Планирование проекта}

\subsection{Понятие планирования}

\begin{definition}
\textbf{Планирование проекта} --- это процесс определения целей, содержания, работ, сроков, ресурсов, бюджета, рисков, коммуникаций и способов контроля, необходимых для успешного выполнения проекта.
\end{definition}

Планирование отвечает на вопросы:
\begin{itemize}
    \item что нужно сделать?
    \item кто будет делать?
    \item когда будет сделано?
    \item сколько это будет стоить?
    \item какие ресурсы нужны?
    \item какие риски существуют?
    \item как будет контролироваться выполнение?
    \item как будут приниматься изменения?
\end{itemize}

\section{План проекта}

\begin{definition}
\textbf{План проекта} --- это согласованный документ или набор документов, определяющих способ выполнения, мониторинга, контроля и завершения проекта.
\end{definition}

План проекта обычно включает:
\begin{enumerate}
    \item цели и обоснование проекта;
    \item содержание проекта;
    \item иерархическую структуру работ;
    \item календарный график;
    \item бюджет;
    \item план ресурсов;
    \item план качества;
    \item план рисков;
    \item план коммуникаций;
    \item план закупок;
    \item план управления изменениями;
    \item критерии приёмки;
    \item процедуры отчётности и контроля.
\end{enumerate}

\subsection{Отличие плана от расписания}

\textbf{План проекта} шире, чем календарный график.

\begin{itemize}
    \item План отвечает на вопрос: \textit{как управлять проектом в целом?}
    \item Расписание отвечает на вопрос: \textit{когда выполняются конкретные работы?}
\end{itemize}

\section{Иерархическая структура работ}

\begin{definition}
\textbf{Иерархическая структура работ}, или \textbf{WBS} --- это декомпозиция результата проекта и необходимых работ на управляемые элементы.
\end{definition}

WBS строится по принципу:
\[
\text{Проект} \rightarrow \text{крупные результаты} \rightarrow \text{подрезультаты} \rightarrow \text{пакеты работ}.
\]

\begin{example}
Проект: система записи студентов на консультации.

\begin{enumerate}
    \item Анализ требований.
    \begin{enumerate}
        \item Интервью со студентами.
        \item Интервью с преподавателями.
        \item Спецификация требований.
    \end{enumerate}
    \item Разработка.
    \begin{enumerate}
        \item Модуль авторизации.
        \item Модуль расписания.
        \item Модуль уведомлений.
    \end{enumerate}
    \item Тестирование.
    \begin{enumerate}
        \item Функциональное тестирование.
        \item Нагрузочное тестирование.
        \item Приёмочное тестирование.
    \end{enumerate}
    \item Ввод в эксплуатацию.
\end{enumerate}
\end{example}

\section{Алгоритм планирования проекта}

\subsection{Шаги алгоритма}

\begin{enumerate}
    \item Определить цели и критерии успеха.
    \item Зафиксировать содержание проекта.
    \item Построить WBS.
    \item Определить перечень работ.
    \item Определить зависимости между работами.
    \item Оценить длительность работ.
    \item Оценить необходимые ресурсы.
    \item Построить календарный график.
    \item Оценить стоимость.
    \item Выявить риски.
    \item Разработать планы реагирования на риски.
    \item Согласовать план с заинтересованными сторонами.
    \item Утвердить базовый план.
\end{enumerate}

\subsection{Иллюстрация}

\begin{center}
\begin{tikzpicture}[
    node distance=1.05cm,
    every node/.style={draw, rounded corners, align=center, minimum width=7.5cm, minimum height=0.7cm},
    arrow/.style={->, thick}
]
\node (a) {Цели и содержание};
\node (b) [below of=a] {WBS};
\node (c) [below of=b] {Работы и зависимости};
\node (d) [below of=c] {Оценки сроков, ресурсов и стоимости};
\node (e) [below of=d] {График и бюджет};
\node (f) [below of=e] {Риски и планы реагирования};
\node (g) [below of=f] {Утверждённый базовый план};

\draw[arrow] (a) -- (b);
\draw[arrow] (b) -- (c);
\draw[arrow] (c) -- (d);
\draw[arrow] (d) -- (e);
\draw[arrow] (e) -- (f);
\draw[arrow] (f) -- (g);
\end{tikzpicture}
\end{center}

\section{Сетевое планирование}

Для индустриальных проектов часто применяют сетевые модели, в которых работы изображаются вершинами или дугами графа, а зависимости между ними --- направленными рёбрами.

\begin{definition}
\textbf{Сетевая модель проекта} --- это ориентированный граф, отражающий порядок выполнения работ и зависимости между ними.
\end{definition}

Если работа $B$ не может начаться до завершения работы $A$, то существует зависимость:
\[
A \rightarrow B.
\]

\subsection{Критический путь}

\begin{definition}
\textbf{Критический путь} --- это самая длинная по длительности цепочка зависимых работ от начала до завершения проекта. Длительность критического пути определяет минимально возможную длительность проекта при заданных зависимостях и длительностях работ.
\end{definition}

Работы на критическом пути имеют нулевой или минимальный резерв времени. Задержка такой работы обычно приводит к задержке всего проекта.

\section{Алгоритм метода критического пути}

\subsection{Входные данные}

Для каждой работы известны:
\begin{itemize}
    \item длительность;
    \item предшествующие работы.
\end{itemize}

\subsection{Основные обозначения}

Для работы $i$:
\begin{itemize}
    \item $d_i$ --- длительность работы;
    \item $ES_i$ --- раннее начало;
    \item $EF_i$ --- раннее окончание;
    \item $LS_i$ --- позднее начало;
    \item $LF_i$ --- позднее окончание;
    \item $R_i$ --- полный резерв.
\end{itemize}

Формулы:
\[
EF_i = ES_i + d_i.
\]

Если у работы $i$ есть предшественники $Pred(i)$, то:
\[
ES_i = \max_{j \in Pred(i)} EF_j.
\]

Если предшественников нет:
\[
ES_i = 0.
\]

Для обратного прохода:
\[
LS_i = LF_i - d_i.
\]

Если у работы $i$ есть последователи $Succ(i)$, то:
\[
LF_i = \min_{j \in Succ(i)} LS_j.
\]

Резерв:
\[
R_i = LS_i - ES_i = LF_i - EF_i.
\]

Работа критична, если:
\[
R_i = 0.
\]

\subsection{Шаги алгоритма}

\begin{enumerate}
    \item Построить ориентированный ациклический граф работ.
    \item Выполнить прямой проход:
    \begin{enumerate}
        \item для начальных работ положить $ES=0$;
        \item вычислить $EF=ES+d$;
        \item для каждой следующей работы взять максимум ранних окончаний предшественников.
    \end{enumerate}
    \item Определить длительность проекта как максимум $EF$ среди завершающих работ.
    \item Выполнить обратный проход:
    \begin{enumerate}
        \item для завершающих работ положить $LF$ равным длительности проекта;
        \item вычислить $LS=LF-d$;
        \item для предшествующих работ взять минимум поздних начал последователей.
    \end{enumerate}
    \item Вычислить резервы $R=LS-ES$.
    \item Найти работы с нулевым резервом.
    \item Цепочка таких работ образует критический путь.
\end{enumerate}

\subsection{Минимальный пример}

Пусть проект состоит из работ:

\begin{center}
\begin{tabular}{|c|c|c|}
\hline
Работа & Длительность & Предшественники \\
\hline
A & 3 & --- \\
\hline
B & 2 & A \\
\hline
C & 4 & A \\
\hline
D & 2 & B, C \\
\hline
\end{tabular}
\end{center}

Граф зависимостей:

\[
A \rightarrow B \rightarrow D,
\qquad
A \rightarrow C \rightarrow D.
\]

Прямой проход:
\[
ES_A=0,\quad EF_A=3.
\]
\[
ES_B=3,\quad EF_B=5.
\]
\[
ES_C=3,\quad EF_C=7.
\]
\[
ES_D=\max(EF_B,EF_C)=\max(5,7)=7,\quad EF_D=9.
\]

Длительность проекта:
\[
T=9.
\]

Обратный проход:
\[
LF_D=9,\quad LS_D=7.
\]
\[
LF_B=LS_D=7,\quad LS_B=5.
\]
\[
LF_C=LS_D=7,\quad LS_C=3.
\]
\[
LF_A=\min(LS_B,LS_C)=\min(5,3)=3,\quad LS_A=0.
\]

Резервы:
\[
R_A=0-0=0,
\]
\[
R_B=5-3=2,
\]
\[
R_C=3-3=0,
\]
\[
R_D=7-7=0.
\]

Критический путь:
\[
A \rightarrow C \rightarrow D.
\]

Длительность критического пути:
\[
3+4+2=9.
\]

\section{Управление содержанием проекта}

\begin{definition}
\textbf{Содержание проекта} --- это совокупность результатов и работ, которые должны быть выполнены для достижения целей проекта.
\end{definition}

Управление содержанием включает:
\begin{enumerate}
    \item сбор требований;
    \item определение границ проекта;
    \item создание WBS;
    \item подтверждение содержания;
    \item контроль изменений содержания.
\end{enumerate}

\subsection{Scope creep}

\begin{definition}
\textbf{Scope creep} --- неконтролируемое расширение содержания проекта без соответствующего пересмотра сроков, бюджета и ресурсов.
\end{definition}

Пример: заказчик просит добавить новые функции ``заодно'', но сроки и бюджет остаются прежними.

Способы предотвращения:
\begin{itemize}
    \item чёткая спецификация требований;
    \item формальная процедура изменения;
    \item матрица трассируемости;
    \item регулярное согласование с заказчиком;
    \item управление ожиданиями заинтересованных сторон.
\end{itemize}

\section{Управление изменениями}

Изменения в проекте неизбежны, особенно если проект длительный или инновационный. Поэтому важна не попытка полностью запретить изменения, а управляемая процедура их рассмотрения.

\subsection{Алгоритм обработки запроса на изменение}

\begin{enumerate}
    \item Получить запрос на изменение.
    \item Зарегистрировать запрос.
    \item Описать причину и суть изменения.
    \item Оценить влияние на содержание, сроки, стоимость, качество и риски.
    \item Подготовить варианты решения:
    \begin{itemize}
        \item принять изменение;
        \item отклонить изменение;
        \item отложить изменение;
        \item заменить альтернативным решением.
    \end{itemize}
    \item Передать запрос на рассмотрение уполномоченной группе или заказчику.
    \item Принять решение.
    \item Обновить план, требования, график и бюджет при необходимости.
    \item Довести решение до заинтересованных сторон.
\end{enumerate}

\section{Управление рисками}

\begin{definition}
\textbf{Риск проекта} --- это неопределённое событие или условие, которое в случае наступления оказывает положительное или отрицательное влияние на цели проекта.
\end{definition}

Обычно под риском чаще понимают угрозу, но в проектном управлении риск может быть и возможностью.

\subsection{Основные этапы управления рисками}

\begin{enumerate}
    \item идентификация рисков;
    \item качественная оценка;
    \item количественная оценка, если необходима;
    \item планирование реагирования;
    \item мониторинг рисков.
\end{enumerate}

\subsection{Матрица вероятности и влияния}

\begin{center}
\begin{tabular}{|c|c|c|c|}
\hline
 & Низкое влияние & Среднее влияние & Высокое влияние \\
\hline
Высокая вероятность & Средний риск & Высокий риск & Критический риск \\
\hline
Средняя вероятность & Низкий риск & Средний риск & Высокий риск \\
\hline
Низкая вероятность & Низкий риск & Низкий риск & Средний риск \\
\hline
\end{tabular}
\end{center}

\subsection{Стратегии реагирования на угрозы}

\begin{itemize}
    \item \textbf{Избежать}: изменить план так, чтобы риск исчез.
    \item \textbf{Снизить}: уменьшить вероятность или влияние риска.
    \item \textbf{Передать}: переложить ответственность на другую сторону, например через страхование или контракт.
    \item \textbf{Принять}: осознанно оставить риск без активных мер, подготовив резерв.
\end{itemize}

\begin{example}
Риск: поставщик задержит оборудование.

Варианты реакции:
\begin{itemize}
    \item выбрать другого поставщика;
    \item заказать оборудование заранее;
    \item включить штрафы в договор;
    \item подготовить временную замену.
\end{itemize}
\end{example}

\section{Управление коммуникациями}

Коммуникации необходимы для координации работ и управления ожиданиями.

План коммуникаций определяет:
\begin{itemize}
    \item кому нужна информация;
    \item какая информация нужна;
    \item как часто она предоставляется;
    \item в каком формате;
    \item кто отвечает за подготовку;
    \item какие каналы используются.
\end{itemize}

Пример фрагмента плана коммуникаций:

\begin{center}
\begin{tabular}{|p{3cm}|p{3cm}|p{3cm}|p{3cm}|}
\hline
Получатель & Информация & Периодичность & Формат \\
\hline
Заказчик & Статус проекта & Еженедельно & Отчёт и встреча \\
\hline
Команда & Задачи и блокеры & Ежедневно & Короткое совещание \\
\hline
Руководство & Бюджет и риски & Ежемесячно & Презентация \\
\hline
\end{tabular}
\end{center}

\section{Контроль исполнения проекта}

План бесполезен без контроля. Контроль проекта включает:
\begin{enumerate}
    \item сбор фактических данных;
    \item сравнение факта с планом;
    \item выявление отклонений;
    \item анализ причин;
    \item корректирующие действия;
    \item обновление прогнозов.
\end{enumerate}

Основные контролируемые параметры:
\begin{itemize}
    \item сроки;
    \item стоимость;
    \item качество;
    \item содержание;
    \item риски;
    \item загрузка ресурсов;
    \item удовлетворённость заказчика.
\end{itemize}

\section{Базовый план}

\begin{definition}
\textbf{Базовый план} --- утверждённая версия плана проекта, используемая как основа для контроля исполнения.
\end{definition}

Обычно выделяют:
\begin{itemize}
    \item базовое содержание;
    \item базовый календарный график;
    \item базовый бюджет.
\end{itemize}

Если фактическое исполнение отличается от базового плана, руководитель проекта анализирует отклонения и принимает управленческие решения.

\section{Документы проекта}

Типовой набор документов:
\begin{enumerate}
    \item бизнес-обоснование;
    \item устав проекта;
    \item реестр заинтересованных сторон;
    \item спецификация требований;
    \item матрица трассируемости требований;
    \item WBS;
    \item календарный график;
    \item бюджет;
    \item реестр рисков;
    \item план коммуникаций;
    \item план качества;
    \item журнал изменений;
    \item отчёты о статусе;
    \item протоколы испытаний;
    \item акты приёмки;
    \item итоговый отчёт.
\end{enumerate}

\section{Роли в проекте}

\begin{description}
    \item[Заказчик] определяет потребность, финансирует или принимает результат.
    \item[Спонсор] обеспечивает организационную поддержку и ресурсы.
    \item[Руководитель проекта] отвечает за достижение целей проекта.
    \item[Команда проекта] выполняет работы по созданию результата.
    \item[Пользователи] используют результат проекта.
    \item[Подрядчики и поставщики] предоставляют внешние работы, услуги или материалы.
    \item[Куратор или управляющий комитет] принимает ключевые управленческие решения.
\end{description}

\section{Пример структуры плана проекта}

\begin{enumerate}
    \item Общие сведения.
    \begin{enumerate}
        \item Название проекта.
        \item Основание для запуска.
        \item Цели и ожидаемая ценность.
    \end{enumerate}

    \item Содержание проекта.
    \begin{enumerate}
        \item Включённые работы.
        \item Исключённые работы.
        \item Основные результаты.
    \end{enumerate}

    \item Требования и критерии приёмки.

    \item Организационная структура.
    \begin{enumerate}
        \item Роли.
        \item Ответственность.
        \item Матрица RACI.
    \end{enumerate}

    \item Иерархическая структура работ.

    \item Календарный план.

    \item Ресурсный план.

    \item Бюджет.

    \item План качества.

    \item План рисков.

    \item План коммуникаций.

    \item План закупок.

    \item План управления изменениями.

    \item Порядок контроля и отчётности.

    \item Порядок закрытия проекта.
\end{enumerate}

\section{Матрица ответственности RACI}

\begin{definition}
\textbf{RACI} --- это матрица распределения ответственности по задачам проекта.
\end{definition}

Обозначения:
\begin{itemize}
    \item $R$ --- Responsible: выполняет работу;
    \item $A$ --- Accountable: несёт итоговую ответственность;
    \item $C$ --- Consulted: консультирует;
    \item $I$ --- Informed: информируется.
\end{itemize}

Пример:

\begin{center}
\begin{tabular}{|p{4cm}|c|c|c|}
\hline
Задача & Руководитель проекта & Аналитик & Заказчик \\
\hline
Сбор требований & A & R & C \\
\hline
Согласование требований & R & C & A \\
\hline
Разработка графика & A/R & C & I \\
\hline
Приёмка результата & R & C & A \\
\hline
\end{tabular}
\end{center}

\section{Пример мини-плана проекта}

Проект: создание системы записи студентов на консультации.

\subsection*{Цель}

Создать веб-систему, позволяющую студентам записываться на консультации преподавателей без ручного согласования по электронной почте.

\subsection*{Содержание}

В проект включены:
\begin{itemize}
    \item авторизация пользователей;
    \item просмотр расписания преподавателей;
    \item запись на свободные слоты;
    \item отмена записи;
    \item уведомления.
\end{itemize}

Не включены:
\begin{itemize}
    \item мобильное приложение;
    \item интеграция с внешними календарями;
    \item автоматическое формирование учебной нагрузки.
\end{itemize}

\subsection*{Основные работы}

\begin{enumerate}
    \item Сбор требований --- 5 дней.
    \item Проектирование интерфейса --- 4 дня.
    \item Разработка backend --- 10 дней.
    \item Разработка frontend --- 10 дней.
    \item Тестирование --- 5 дней.
    \item Внедрение --- 2 дня.
\end{enumerate}

\subsection*{Риски}

\begin{center}
\begin{tabular}{|p{5cm}|c|c|p{5cm}|}
\hline
Риск & Вероятность & Влияние & Реакция \\
\hline
Преподаватели не согласуют требования вовремя & Средняя & Высокое & Назначить ответственного и сроки согласования \\
\hline
Недостаточная производительность & Низкая & Среднее & Провести нагрузочное тестирование \\
\hline
Изменение требований во время разработки & Высокая & Среднее & Ввести процедуру контроля изменений \\
\hline
\end{tabular}
\end{center}

\section{Итоговые положения}

Управление проектом объединяет техническую, организационную и коммуникационную деятельность. Проект отличается временным характером, уникальностью результата и наличием ограничений. Творческие проекты требуют гибкости и допускают существенное уточнение результата в ходе работы. Индустриальные проекты чаще строятся по фазовой модели, требуют формальной документации, контроля качества, управления рисками и строгого согласования изменений.

Сбор требований является основой для определения содержания проекта. Неполные или противоречивые требования приводят к ошибкам планирования, росту стоимости и задержкам. Планирование превращает цели и требования в управляемую структуру работ, график, бюджет, распределение ответственности и систему контроля. Хороший план не устраняет неопределённость полностью, но делает проект прозрачным и управляемым.

\end{document}
